\documentclass[11pt]{letter}
\usepackage[utf8]{inputenc}
\usepackage[T1]{fontenc}
\usepackage{lmodern}
\usepackage[margin=1.1in]{geometry}
\usepackage{amsmath,amssymb}
\usepackage{microtype}
\usepackage{hyperref}
\hypersetup{hidelinks}

\address{David H.\ Silver\\
Independent researcher\\
Poughkeepsie, NY 12603, USA\\
\href{mailto:silverdavi@gmail.com}{silverdavi@gmail.com}\\
ORCID: \href{https://orcid.org/0000-0002-3071-304X}{0000-0002-3071-304X}}
\signature{David H.\ Silver}
\begin{document}
\begin{letter}{Editor-in-Chief\\ \emph{Journal of Peace Research}\\ Oxford University Press / PRIO}

\opening{Dear Editor,}

I am submitting the enclosed manuscript, \emph{``Bounds and Estimators for the Civilian-to-Combatant Death Ratio from Aggregation of Conflicting Sources,''} for consideration as a Regular Article in the \emph{Journal of Peace Research}.

Conflict-data research has long acknowledged that casualty figures are produced by sources with structural reasons to bias their estimates---a tension visible in every UCDP, ACLED, IBC, and ministry-of-health series (Sundberg and Melander, 2013; Raleigh et al., 2010; Lacina and Gleditsch, 2005), and made acute by the recent UCDP-coder reassessment of Vesco et al.\ (2026). Within the conflict-research tradition, the natural response has been either to defer to a single channel or to publish wide ``low--high'' bands. The paper proposes a third path: treat the civilian-to-combatant ratio as a \emph{partially identified} parameter and identify the polytope of values jointly consistent with all sources, before any single channel is trusted.

The paper's four contributions are: (i) a sharp identified set for the combatant share $q=D_M/D$ under bounded differential exposure of civilian adult men; (ii) a precision-weighted aggregation rule for multiple demographic sources; (iii) a \emph{contradiction radius} $\rho$ that scores any disputed claim by the standardised perturbation of independent inputs needed to make it consistent; and (iv) a closed-form bias-robust credible interval under Huber $\varepsilon$-contamination of sources. The framework is symmetric between belligerents: it rejects ``too many combatants'' and ``zero combatants'' against the same feasible set.

Applied to the Gaza war 2023--2026, the framework finds that the public Israel Defense Forces claim of $17$--$25$k combatants killed has $\rho\gtrsim 20$ standard errors. A spatial Bayesian model on the same independently-measured inputs (PCBS demographics, OHCHR identifications, Gaza Ministry of Health records, OCHA totals, IISS/RUSI manpower) gives $q=2.5\%$ with 95\% credible interval $[1.6\%,3.8\%]$ and an implied civilian-to-combatant ratio of approximately $43{:}1$. The diagnostic is then run on 81 conflicts from 1900--2026, reporting which admit a sharp diagnostic (Iraq 2003--11, Gaza 2023--26, Tigray, Yemen) and which do not (most of the 20th century).

The paper fits naturally into JPR's tradition of methodologically rigorous conflict-data work: Eck and Hultman (2007) on one-sided violence data, Sundberg and Melander (2013) on UCDP-GED, Raleigh et al.\ (2010) on ACLED, and the recent VIEWS Prediction Challenge contributions in vol.~62 (2025). It is complementary to (rather than competing with) capture--recapture work in human-rights statistics (Lum, Price and Banks, 2013; Khatib et al., 2025) and demographic-uncertainty Bayesian work in population health (G\'{o}mez-Ugarte et al., 2025; Alburez-Gutierrez et al., 2024).

The manuscript runs approximately $7{,}500$ words plus an online supplement containing the per-conflict tables for the 81-conflict survey and full proofs of regularity conditions. Code, simulator, and Gaza inputs will be deposited in the JPR Replication Data archive on acceptance, in compliance with the journal's Data Replication Policy. This work has not been considered or submitted elsewhere, and I declare no conflicts of interest.

\textbf{About the author.} I am an independent researcher with peer-reviewed publications in \emph{Nature}, \emph{PNAS}, \emph{IEEE Trans.\ Pattern Analysis and Machine Intelligence}, \emph{Human Reproduction}, and \emph{PLOS ONE}, and I held a Microsoft Research PhD Fellowship in Biomathematics, Bioinformatics, and Computational Biology jointly between the Technion and Microsoft Research Cambridge (UK). I am the author of the peer-reviewed monograph \emph{Beyond Popular Science} (Open Book Publishers, Cambridge, 2026; DOI 10.11647/OBP.0526). The present manuscript is conducted independently and is unaffiliated with any of my industry roles.

Suggested reviewers include Michael Spagat (Royal Holloway), Patrick Ball (HRDAG), Lisa Hultman (Uppsala), Diego Alburez-Gutierrez (MPIDR), Erik Melander (Uppsala/UCDP), and H{\aa}vard Hegre (Uppsala/PRIO).

\closing{Sincerely,}

\end{letter}
\end{document}
